% Template for PLoS
% Version 3.4 January 2017
%
% % % % % % % % % % % % % % % % % % % % % %
%
% -- IMPORTANT NOTE
%
% This template contains comments intended 
% to minimize problems and delays during our production 
% process. Please follow the template instructions
% whenever possible.
%
%
% -- FIGURES AND TABLES
%
% Please include tables/figure captions directly after the paragraph where they are first cited in the text.
%
% DO NOT INCLUDE GRAPHICS IN YOUR MANUSCRIPT
% - Figures should be uploaded separately from your manuscript file. 
% - Figures generated using LaTeX should be extracted and removed from the PDF before submission. 
% - Figures containing multiple panels/subfigures must be combined into one image file before submission.
% For figure citations, please use "Fig" instead of "Figure".
% See http://journals.plos.org/plosone/s/figures for PLOS figure guidelines.
%
% Tables should be cell-based and may not contain:
% - spacing/line breaks within cells to alter layout or alignment
% - do not nest tabular environments (no tabular environments within tabular environments)
% - no graphics or colored text (cell background color/shading OK)
% See http://journals.plos.org/plosone/s/tables for table guidelines.
%
% For tables that exceed the width of the text column, use the adjustwidth environment as illustrated in the example table in text below.
%
% % % % % % % % % % % % % % % % % % % % % % % %
%
% -- EQUATIONS, MATH SYMBOLS, SUBSCRIPTS, AND SUPERSCRIPTS
%
% IMPORTANT
% Below are a few tips to help format your equations and other special characters according to our specifications. For more tips to help reduce the possibility of formatting errors during conversion, please see our LaTeX guidelines at http://journals.plos.org/plosone/s/latex
%
% For inline equations, please be sure to include all portions of an equation in the math environment.  For example, x$^2$ is incorrect; this should be formatted as $x^2$ (or $\mathrm{x}^2$ if the romanized font is desired).
%
% Do not include text that is not math in the math environment. For example, CO2 should be written as CO\textsubscript{2} instead of CO$_2$.
%
% Please add line breaks to long display equations when possible in order to fit size of the column. 
%
% For inline equations, please do not include punctuation (commas, etc) within the math environment unless this is part of the equation.
%
% When adding superscript or subscripts outside of brackets/braces, please group using {}.  For example, change "[U(D,E,\gamma)]^2" to "{[U(D,E,\gamma)]}^2". 
%
% Do not use \cal for caligraphic font.  Instead, use \mathcal{}
%
% % % % % % % % % % % % % % % % % % % % % % % % 
%
% Please contact latex@plos.org with any questions.
%
% % % % % % % % % % % % % % % % % % % % % % % %

\documentclass[10pt,letterpaper]{article}
\usepackage[top=0.85in,left=2.75in,footskip=0.75in]{geometry}

% amsmath and amssymb packages, useful for mathematical formulas and symbols
\usepackage{amsmath,amssymb}

% Use adjustwidth environment to exceed column width (see example table in text)
\usepackage{changepage}

% Use Unicode characters when possible
\usepackage[utf8x]{inputenc}

% textcomp package and marvosym package for additional characters
\usepackage{textcomp,marvosym}

% cite package, to clean up citations in the main text. Do not remove.
\usepackage{cite}

% Use nameref to cite supporting information files (see Supporting Information section for more info)
\usepackage{nameref,hyperref}

% line numbers
\usepackage[right]{lineno}

% ligatures disabled
\usepackage{microtype}
\DisableLigatures[f]{encoding = *, family = * }

% color can be used to apply background shading to table cells only
\usepackage[table]{xcolor}

% array package and thick rules for tables
\usepackage{array}

% create "+" rule type for thick vertical lines
\newcolumntype{+}{!{\vrule width 2pt}}

% create \thickcline for thick horizontal lines of variable length
\newlength\savedwidth
\newcommand\thickcline[1]{%
	\noalign{\global\savedwidth\arrayrulewidth\global\arrayrulewidth 2pt}%
	\cline{#1}%
	\noalign{\vskip\arrayrulewidth}%
	\noalign{\global\arrayrulewidth\savedwidth}%
}

% \thickhline command for thick horizontal lines that span the table
\newcommand\thickhline{\noalign{\global\savedwidth\arrayrulewidth\global\arrayrulewidth 2pt}%
	\hline
	\noalign{\global\arrayrulewidth\savedwidth}}


% Remove comment for double spacing
%\usepackage{setspace} 
%\doublespacing

% Text layout
\raggedright
\setlength{\parindent}{0.5cm}
\textwidth 5.25in 
\textheight 8.75in

% Bold the 'Figure #' in the caption and separate it from the title/caption with a period
% Captions will be left justified
\usepackage[aboveskip=1pt,labelfont=bf,labelsep=period,justification=raggedright,singlelinecheck=off]{caption}
\renewcommand{\figurename}{Fig}

% Use the PLoS provided BiBTeX style
\bibliographystyle{plos2015}

% Remove brackets from numbering in List of References
\makeatletter
\renewcommand{\@biblabel}[1]{\quad#1.}
\makeatother

% Leave date blank
\date{}

% Header and Footer with logo
\usepackage{lastpage,fancyhdr,graphicx}
\usepackage{epstopdf}
\pagestyle{myheadings}
\pagestyle{fancy}
\fancyhf{}
\setlength{\headheight}{27.023pt}
\lhead{\includegraphics[width=2.0in]{PLOS-submission.eps}}
\rfoot{\thepage/\pageref{LastPage}}
\renewcommand{\footrule}{\hrule height 2pt \vspace{2mm}}
\fancyheadoffset[L]{2.25in}
\fancyfootoffset[L]{2.25in}
\lfoot{\sf PLOS}

%% Include all macros below
\DeclareMathOperator*{\argmin}{arg\,min}
\DeclareMathOperator*{\argmax}{arg\,max}
%% END MACROS SECTION


\begin{document}
	\vspace*{0.2in}
	
	% Title must be 250 characters or less.
	\begin{flushleft}
		{\Large
			\textbf\newline{Predicting signal peptides using machine learning} % Please use "sentence case" for title and headings (capitalize only the first word in a title (or heading), the first word in a subtitle (or subheading), and any proper nouns).
		}
		\newline
		% Insert author names, affiliations and corresponding author email (do not include titles, positions, or degrees).
		\\
		Jonas Andersson
		\\
		\bigskip
		Computer Science and Communications, KTH Royal Institute of Technology, Stockholm, Sweden
		\bigskip
		
		% Insert additional author notes using the symbols described below. Insert symbol callouts after author names as necessary.
		% 
		% Remove or comment out the author notes below if they aren't used.
		%
		% Primary Equal Contribution Note
		%\Yinyang These authors contributed equally to this work.
		
		% Additional Equal Contribution Note
		% Also use this double-dagger symbol for special authorship notes, such as senior authorship.
		%\ddag These authors also contributed equally to this work.
		
		% Current address notes
		%\textcurrency Current Address: Dept/Program/Center, Institution Name, City, State, Country % change symbol to "\textcurrency a" if more than one current address note
		% \textcurrency b Insert second current address 
		% \textcurrency c Insert third current address
		
		% Deceased author note
		%\dag Deceased
		
		% Group/Consortium Author Note
		%\textpilcrow Membership list can be found in the Acknowledgments section.
		
		% Use the asterisk to denote corresponding authorship and provide email address in note below.
		* jonand8@kth.se
		
	\end{flushleft}
	% Please keep the abstract below 300 words
	\section*{Abstract}
	Abstract about using machine learning to predict signal peptides
	
	
	% Please keep the Author Summary between 150 and 200 words
	% Use first person. PLOS ONE authors please skip this step. 
	% Author Summary not valid for PLOS ONE submissions.   
	%\section*{Author summary}
	%Lorem ipsum dolor sit amet, consectetur adipiscing elit. Curabitur eget %porta erat. Morbi consectetur est vel gravida pretium. Suspendisse ut dui %eu ante cursus gravida non sed sem. Nullam sapien tellus, commodo id velit %id, eleifend volutpat quam. Phasellus mauris velit, dapibus finibus %elementum vel, pulvinar non tellus. Nunc pellentesque pretium diam, quis %maximus dolor faucibus id. Nunc convallis sodales ante, ut ullamcorper est %egestas vitae. Nam sit amet enim ultrices, ultrices elit pulvinar, %volutpat risus.
	
	\linenumbers
	
	% Use "Eq" instead of "Equation" for equation citations.
	\section*{Introduction}
	The localization of proteins is mostly determined by a biological trafficking system. When proteins are synthesized on cytoplasmic ribosomes and have several possible destinations such as the nuclear, mitochondria, or peroxisomes. Signal peptides direct proteins to their proper destination after synthesis, keeping them from arriving at the wrong cellular location. Signal peptides, short chains of amino acids, act as postal codes for their destination. These sorting signals, a N-terminal signal peptide, are removed or cleaved off. If the protein does not have this signal peptide then it remains in the cytoplasm. The structure of the signal peptide consists of a positively charged N-region (1-5 residues long), a hydrophobic h-region (7-15 residues long), and finally a polar C-region (3-7 residues long). The cleavage point is in the C-region, but the location varies between proteins \cite{thesis}. 
	
	There exist a number of computational methods for predicting these sorting signals using distinctly different approaches. It is useful to locate these signals as it can potentially add insight into protein translocation, how chromosomes are rearranged. Some examples of application domains for signal prediction include drug discovery and gene therapy \cite{thesis}. Numbering among these computational approaches are neural networks(NNs), support vector machines(SVMs), hidden markov models (HMMs), decision trees, and ensemble learning methods. These algorithms are encompassed in field of computer science called machine learning. Machine learning (ML) algorithms learn and make decisions from data \cite{wikipedia}. ML comprises a broad spectrum of approaches; much of the work done on signal prediction is approached as a classification problem. In classification, a probability distribution is drawn over the classes, and we choose the class that has the highest probability \cite{jurafsky}. $$ \hat{c} = \argmax_{c \in C}{ \Pr(c \vert x)} $$
	
	A sequence classifier is a model whose job is to assign a label to some part of a sequence \cite{jurafsky}. 
	
	The advantage of HMM's is that they can be constructed from only positive examples.
	
	
	\section*{Materials and methods}
	\subsection*{Etiam eget sapien nibh}
	
	% For figure citations, please use "Fig" instead of "Figure".
	
	% Place figure captions after the first paragraph in which they are cited.
	\begin{figure}[!h]
		\caption{{\bf Bold the figure title.}
			Figure caption text here, please use this space for the figure panel descriptions instead of using subfigure commands. A: Lorem ipsum dolor sit amet. B: Consectetur adipiscing elit.}
		\label{fig1}
	\end{figure}
	
	% Results and Discussion can be combined.
	\section*{Results}
	
	% Place tables after the first paragraph in which they are cited.
	\begin{table}[!ht]
		\begin{adjustwidth}{-2.25in}{0in} % Comment out/remove adjustwidth environment if table fits in text column.
			\centering
			\caption{
				{\bf Table caption Nulla mi mi, venenatis sed ipsum varius, volutpat euismod diam.}}
			\begin{tabular}{|l+l|l|l|l|l|l|l|}
				\hline
				\multicolumn{4}{|l|}{\bf Heading1} & \multicolumn{4}{|l|}{\bf Heading2}\\ \thickhline
				$cell1 row1$ & cell2 row 1 & cell3 row 1 & cell4 row 1 & cell5 row 1 & cell6 row 1 & cell7 row 1 & cell8 row 1\\ \hline
				$cell1 row2$ & cell2 row 2 & cell3 row 2 & cell4 row 2 & cell5 row 2 & cell6 row 2 & cell7 row 2 & cell8 row 2\\ \hline
				$cell1 row3$ & cell2 row 3 & cell3 row 3 & cell4 row 3 & cell5 row 3 & cell6 row 3 & cell7 row 3 & cell8 row 3\\ \hline
			\end{tabular}
			\begin{flushleft} Table notes Phasellus venenatis, tortor nec vestibulum mattis, massa tortor interdum felis, nec pellentesque metus tortor nec nisl. Ut ornare mauris tellus, vel dapibus arcu suscipit sed.
			\end{flushleft}
			\label{table1}
		\end{adjustwidth}
	\end{table}
	
	
	%PLOS does not support heading levels beyond the 3rd (no 4th level headings).
	\subsection*{nunc blandit a tortor}
	\subsubsection*{3rd level heading} 
	
	\begin{enumerate}
		\item{react}
		\item{diffuse free particles}
		\item{increment time by dt and go to 1}
	\end{enumerate}
	
	
	\subsection*{Sed ac quam id nisi malesuada congue}

	
	\begin{itemize}
		\item First bulleted item.
		\item Second bulleted item.
		\item Third bulleted item.
	\end{itemize}
	
	\section*{Discussion}
	
	\section*{Conclusion}
	
	

	\nolinenumbers
	
	% Either type in your references using
	% \begin{thebibliography}{}
	% \bibitem{}
	% Text
	% \end{thebibliography}
	%
	% or
	%
	% Compile your BiBTeX database using our plos2015.bst
	% style file and paste the contents of your .bbl file
	% here. See http://journals.plos.org/plosone/s/latex for 
	% step-by-step instructions.
	% 
	\begin{thebibliography}{10}
		
		\bibitem{thesis}
		Ojefua EE.
		\newblock {{F}inding Secretion Signals in Protein Sequences}.
		\newblock Masters thesis. University of Tampere; 2009.
		
		\bibitem{jurafsky}
		Jurafsky D. Martin HM.
		\newblock {{S}peech and language processing: An introduction to natural language processing, computational linguistics, and speech recognition.}
		\newblock {London: Pearson Education; 2009}
		
	\end{thebibliography}
	
	
	
\end{document}

